\documentclass[11pt]{article}
\usepackage[margin=1in]{geometry}
\usepackage{amsmath,amssymb,mathtools}
\usepackage{microtype}
\setlength{\parindent}{0pt}
\setlength{\parskip}{0.7em}

\newcommand{\I}{\mathbb{I}}

\begin{document}

\begin{center}
  {\Large\bfseries Problem 1 -- Question 3}\\[0.3em]
  {\large Implementing \(R_{\mathbf n}(\theta)\) with \(R_Y\) and \(R_Z\)}
\end{center}

Parameterize the unit vector \(\mathbf n\) by polar and azimuthal angles:
\[
 \mathbf n
 =(\sin\lambda\cos\phi,\,
   \sin\lambda\sin\phi,\,
   \cos\lambda).
\]
Let
\[
 W=R_Z(\phi)R_Y(\lambda).
\]
Rotations conjugate Pauli vectors in the same way that their corresponding
three-dimensional rotations act on ordinary vectors. In particular,
\begin{align*}
 WZW^\dagger
 &=R_Z(\phi)
   \bigl(R_Y(\lambda)ZR_Y(-\lambda)\bigr)
   R_Z(-\phi)\\
 &=\sin\lambda\cos\phi\,X
  +\sin\lambda\sin\phi\,Y
  +\cos\lambda\,Z\\
 &=\mathbf n\cdot\boldsymbol{\sigma}.
\end{align*}
Therefore
\begin{align*}
 R_{\mathbf n}(\theta)
 &=\exp\left(-i\frac{\theta}{2}WZW^\dagger\right)\\
 &=W\exp\left(-i\frac{\theta}{2}Z\right)W^\dagger\\
 &=R_Z(\phi)R_Y(\lambda)R_Z(\theta)
   R_Y(-\lambda)R_Z(-\phi).
\end{align*}
Thus a valid gate sequence, listed in chronological order, is
\[
\boxed{
R_Z(-\phi)\;,\quad
R_Y(-\lambda)\;,\quad
R_Z(\theta)\;,\quad
R_Y(\lambda)\;,\quad
R_Z(\phi).
}
\]
The rightmost factor in the matrix product acts first, which explains the
reversal between the product and the chronological list.

\paragraph{Special cases.}
If \(\mathbf n=\hat z\), take \(\lambda=0\), and the construction reduces
to \(R_Z(\theta)\). If \(\mathbf n=-\hat z\), take \(\lambda=\pi\).
At either pole, \(\phi\) is arbitrary and may be set to zero.

\paragraph{Resource count.}
The direct construction uses five elementary rotations. Adjacent rotations
about the same axis can be merged with neighboring circuit gates. In a
larger circuit this often reduces the effective count; for example,
\(R_Z(a)R_Z(b)=R_Z(a+b)\).

\end{document}
